\documentclass{stex}

\libusepackage{naproche}
\libusepackage{copyright}
\libinput{preamble}

\begin{document}
\begin{smodule}{first-order-semantics}

\importmodule[libraries/meta]{preliminaries.ftl}

\begin{sfragment}[id=sec:first-order-semantics]{First-Order Semantics}
  \symdecl*{formula}
  \symdecl*{term}
  \symdef{folsem}[name=first-order semantics,args=1,op=[\![\cdot]\!]]{\comp{[\![}#1\comp{]\!]}}
  %
  \begin{sparagraph}[style=symdoc,for={first-order semantics,folsem}]
    \vardef{Evar}{E}
    %
    \Naproche translates \ForTheL expressions to (a variant of)
    \sr{formula}{first-order formulas} which we refer to as the
    \definame{first-order semantics} of \ForTheL.
    We denote the first-order translation of an expression $\Evar$
    by $\defnotation\folsem\Evar$.
  \end{sparagraph}

  \begin{sfragment}[id=sec:tags]{Tags}
    \symdecl*{tag}
    \symdef{DigTag}{\langle\textsf{dig}\rangle}
    \symdef{DigMultiSubjectTag}{\langle\textsf{dig multi-subject}\rangle}
    \symdef{DigMultiPairwiseTag}{\langle\textsf{dig multi-pairwise}\rangle}
    \symdef{HeadTermTag}{\langle\textsf{head term}\rangle}
    \symdef{InductionHypothesisTag}{\langle\textsf{induction hypothesis}\rangle}
    \symdef{CaseHypothesisTag}{\langle\textsf{case hypothesis}\rangle}
    \symdef{EqualityChainTag}{\langle\textsf{equality chain}\rangle}
    \symdef{GenericMarkTag}{\langle\textsf{generic mark}\rangle}
    \symdef{EvaluationTag}{\langle\textsf{evaluation}\rangle}
    \symdef{ConditionTag}{\langle\textsf{condition}\rangle}
    \symdef{DefinedTag}{\langle\textsf{defined}\rangle}
    \symdef{DomainTag}{\langle\textsf{domain}\rangle}
    \symdef{ReplacementTag}{\langle\textsf{replacement}\rangle}
    \symdef{DomainTaskTag}{\langle\textsf{domain task}\rangle}
    \symdef{ExistenceTaskTag}{\langle\textsf{existence task}\rangle}
    \symdef{UniquenessTaskTag}{\langle\textsf{uniqueness task}\rangle}
    \symdef{ChoiceTaskTag}{\langle\textsf{choice task}\rangle}
    %
    \begin{sparagraph}[style=symdoc,for={tag,DigTag,DigMultiSubjectTag,DigMultiPairwiseTag,HeadTermTag,InductionHypothesisTag,
      CaseHypothesisTag,EqualityChainTag,GenericMarkTag,EvaluationTag,ConditionTag,DefinedTag,DomainTag,ReplacementTag,
      DomainTaskTag,ExistenceTaskTag,UniquenessTaskTag,ChoiceTaskTag}]
      \Sns{formula} can be annotated with one of the following \definame[post=s]{tag}:
      \begin{enumerate}
        \item $\defnotation\DigTag$
        \item $\defnotation\DigMultiSubjectTag$
        \item $\defnotation\DigMultiPairwiseTag$
        \item $\defnotation\HeadTermTag$
        \item $\defnotation\InductionHypothesisTag$
        \item $\defnotation\CaseHypothesisTag$
        \item $\defnotation\EqualityChainTag$
        \item $\defnotation\GenericMarkTag$
        \item $\defnotation\EvaluationTag$
        \item $\defnotation\ConditionTag$
        \item $\defnotation\DefinedTag$
        \item $\defnotation\DomainTag$
        \item $\defnotation\ReplacementTag$
        \item $\defnotation\DomainTaskTag$
        \item $\defnotation\ExistenceTaskTag$
        \item $\defnotation\UniquenessTaskTag$
        \item $\defnotation\ChoiceTaskTag$
      \end{enumerate}
    \end{sparagraph}

    $\GenericMarkTag$, $\EvaluationTag$, $\ConditionTag$, $\DefinedTag$, $\DomainTag$ and $\ReplacementTag$ \sns{tag}
    are used to mark certain parts of map definitions.
    $\DomainTaskTag$, $\ExistenceTaskTag$, $\UniquenessTaskTag$ and $\ChoiceTaskTag$ \sns{tag} are used to mark
    certain parts of well-definedness proof tasks for maps.
  \end{sfragment}

  \begin{sfragment}[id=sec:variables]{Variables}
    \symdecl*{variable name}
    \symdecl*{variable declaration}
    \symdef{constVar}[args=1,name=constant variable]{\maincomp{\textsf{x}}\dobrackets{#1}}
    \symdef{holeVar}[args=1,name=hole variable]{\maincomp{\textsf{?}}\dobrackets{#1}}
    \symdef{slotVar}[name=slot variable]{\textsf{!}}
    \symdef{uVar}[args=1,name=u-variable]{\maincomp{\textsf{u}}\dobrackets{#1}}
    \symdef{hiddenVar}[args=1,name=hidden variable]{\maincomp{\textsf{h}}\dobrackets{#1}}
    \symdef{assumeVar}[args=1,name=assume variable]{\maincomp{\textsf{i}}\dobrackets{#1}}
    \symdef{skolemVar}[args=1,name=Skolem variable]{\maincomp{\textsf{o}}\dobrackets{#1}}
    \symdef{taskVar}[args=1,name=task variable]{\maincomp{\textsf{c}}\dobrackets{#1}}
    \symdef{zVar}[args=1,name=z-variable]{\maincomp{\textsf{z}}\dobrackets{#1}}
    \symdef{wVar}[args=1,name=w-variable]{\maincomp{\textsf{w}}\dobrackets{#1}}
    \symdef{emptyVar}[name=empty variable]{\epsilon_{\textsf{v}}}
    \symdef{defaultVar}[args=1,name=default variable]{\maincomp{\textsf d}\dobrackets{#1}}
    %
    \begin{sparagraph}[style=symdoc,for={variable name,variable declaration,constVar,holeVar,slotVar,uVar,hiddenVar,
      assumeVar,skolemVar,taskVar,zVar,wVar,emptyVar,defaultVar}]
      \vardef{svar}{s}
      \vardef{nvar}{n}
      \vardef{vvar}{v}
      %
      \Naproche knows two kinds of variables: \Definame[post=s]{variable name} and
      \definame[post=s]{variable declaration}.
      The latter is just a \sn{variable name} equipped with a natural number.
      \Sns{variable name} subdivide into the following kinds:
      \begin{enumerate}
        \item \Definame[post=s]{constant variable}, denoted by $\defnotation\constVar\svar$,
          where $\svar$ is a string;
        \item \Definame[post=s]{hole variable}, denoted by $\defnotation\holeVar\svar$,
          where $\svar$ is a string;
        \item \Definame[post=s]{slot variable}, denoted by $\defnotation\slotVar$;
        \item \Definame[post=s]{u-variable}, denoted by $\defnotation\uVar\svar$,
          where $\svar$ is a string;
        \item \Definame[post=s]{hidden variable}, denoted by $\defnotation\hiddenVar\nvar$,
          where $\nvar$ is an integer;
        \item \Definame[post=s]{assume variable}, denoted by $\defnotation\assumeVar\nvar$,
          where $\nvar$ is an integer;
        \item \Definame[post=s]{Skolem variable}, denoted by $\defnotation\skolemVar\nvar$,
          where $\nvar$ is an integer;
        \item \Definame[post=s]{task variable}, denoted by $\defnotation\taskVar\vvar$,
          where $\vvar$ is a \sn{variable name};
        \item \Definame[post=s]{z-variable}, denoted by $\defnotation\zVar\svar$,
          where $\svar$ is a string;
        \item \Definame[post=s]{w-variable}, denoted by $\defnotation\wVar\svar$,
          where $\svar$ is a string;
        \item \Definame[post=s]{empty variable}, denoted by $\defnotation\emptyVar$;
        \item \Definame[post=s]{default variable}, denoted by $\defnotation\defaultVar\svar$,
          where $\svar$ is a string.
        \end{enumerate}
    \end{sparagraph}
  \end{sfragment}

  \begin{sfragment}[id=sec:terms]{Terms}
    \symdecl*{term name}
    \symdef{nameTerm}[args=1,name=name term]{\maincomp{\textsf{n}}\dobrackets{#1}}
    \symdef{symbolicTerm}[args=1,name=symbolic term]{\maincomp{\textsf{s}}\dobrackets{#1}}
    \symdef{notionTerm}[args=1,name=notion term]{\maincomp{\textsf{a}}\dobrackets{#1}}
    \symdef{theTerm}[args=1,name=the-term]{\maincomp{\textsf{the}}\dobrackets{#1}}
    \symdef{uAdjectiveTerm}[args=1,name=unary adjective term]{\maincomp{\textsf{is}}\dobrackets{#1}}
    \symdef{mAdjectiveTerm}[args=1,name=multi adjective term]{\maincomp{\textsf{mis}}\dobrackets{#1}}
    \symdef{uVerbTerm}[args=1,name=unary verb term]{\maincomp{\textsf{do}}\dobrackets{#1}}
    \symdef{mVerbTerm}[args=1,name=multi verb term]{\maincomp{\textsf{mdo}}\dobrackets{#1}}
    \symdef{taskTerm}[args=1,name=task term]{\maincomp{\textsf{tsk}}\dobrackets{#1}}
    \symdef{equalityTerm}[name=equality term]{\textsf{eq}}
    \symdef{ilessTerm}[name=inductively-less term]{\textsf{iless}}
    \symdef{thesisTerm}[name=thesis term]{\textsf{\#TH\#}}
    \symdef{emptyTerm}[name=empty term]{\epsilon_{\textsf{t}}}
    %
    \begin{sparagraph}[style=symdoc,for={term name,nameTerm,symbolicTerm,theTerm,uAdjectiveTerm,mAdjectiveTerm,
    uVerbTerm,mVerbTerm,taskTerm,equalityTerm,ilessTerm,thesisTerm,emptyTerm}]
      \vardef{svar}{s}
      \vardef{nvar}{n}
      %
      \Sns{formula} can consist of \sns{term} which are constructed via a \definame{term name} to
      which a list of \sns{formula} is passed as arguments.
      \Sns{term name} can be of the following forms:
      \begin{enumerate}
        \item \Definame[post=s]{name term}, denoted by $\defnotation\nameTerm\svar$,
          where $\svar$ is a string;
        \item \Definame[post=s]{symbolic term}, denoted by $\defnotation\symbolicTerm\svar$,
          where $\svar$ is a string;
        \item \Definame[post=s]{notion term}, denoted by $\defnotation\notionTerm\svar$,
          where $\svar$ is a string;
        \item \Definame[post=s]{the-term}, denoted by $\defnotation\theTerm\svar$,
          where $\svar$ is a string;
        \item \Definame[post=s]{unary adjective term}, denoted by $\defnotation\uAdjectiveTerm\svar$,
          where $\svar$ is a string;
        \item \Definame[post=s]{multi adjective term}, denoted by $\defnotation\mAdjectiveTerm\svar$,
          where $\svar$ is a string;
        \item \Definame[post=s]{unary verb term}, denoted by $\defnotation\uVerbTerm\svar$,
          where $\svar$ is a string;
        \item \Definame[post=s]{multi verb term}, denoted by $\defnotation\mVerbTerm\svar$,
          where $\svar$ is a string;
        \item \Definame[post=s]{task term}, denoted by $\defnotation\taskTerm\nvar$,
          where $\nvar$ is an integer;
        \item \Definame[post=s]{equality term}, denoted by $\defnotation\equalityTerm$;
        \item \Definame[post=s]{inductively-less term}, denoted by $\defnotation\ilessTerm$;
        \item \Definame[post=s]{thesis term}, denoted by $\defnotation\thesisTerm$;
        \item \Definame[post=s]{empty term}, denoted by $\defnotation\emptyTerm$.
      \end{enumerate}
    \end{sparagraph}

    By default, \Naproche knows the following \sns{notion term}:
    $\notionTerm{\texttt{function}}$,
    $\notionTerm{\texttt{map}}$,
    $\notionTerm{\texttt{set}}$,
    $\notionTerm{\texttt{class}}$,
    $\notionTerm{\texttt{object}}$ and
    $\notionTerm{\texttt{elementOf}}$;
    and the following \sns{name term}:
    $\nameTerm{\texttt{.(.)}}$ (which denotes the map application operation),
    $\nameTerm{\texttt{Dom(.)}}$ (which denotes the domain operation) and
    $\nameTerm{\texttt{(.,.)}}$ (which denotes the ordered pair constructor).
  \end{sfragment}

  \begin{sfragment}[id=sec:formulas]{Formulas}
    \symdef{foltrue}{\top}
    \symdef{folfalse}{\bot}
    \symdef{folthis}{\textsf{this}}
    \symdef{folneg}[args=1]{\mathop{\maincomp\neg}#1}
    \symdef{folconj}[args=a]{\argsep{#1}{\mathbin{\maincomp\wedge}}}
    \symdef{foldisj}[args=a]{\argsep{#1}{\mathbin{\maincomp\vee}}}
    \symdef{folimpl}[args=2]{#1\mathbin{\maincomp\Rightarrow}#2}
    \symdef{folequiv}[args=a]{\argsep{#1}{\mathbin{\maincomp\Leftrightarrow}}}
    \symdef{folforall}[args=2]{\maincomp\forall#1\comp.\,#2}
    \symdef{folexists}[args=2]{\maincomp\exists#1\comp.\,#2}
    \symdef{foltag}[args=2]{#1\mathrel{\maincomp{::}}#2}
    \symdef{foldbindex}[args=1,name=De Brujin index]{#1}
    %
    \begin{sparagraph}[style=symdoc,for={formula,foltrue,folfalse,folneg,folconj,foldisj,folimpl,folequiv,folforall,folexists,foltag,foldbindex,folsem}]
      \vardef{phivar}{\varphi}
      \vardef{psivar}{\psi}
      \vardef{Tvar}{T}
      \vardef{ivar}{i}
      \vardef{xvar}{x}
      \vardef{vvar}{v}
      \vardef{tvar}[args=a]{\maincomp t\dobrackets{#1}}
      \vardef{nvar}{n}
      \varseq{phiseq}{1,\ellipses,\nvar}{\maincomp\varphi_{#1}}
      %
      \Definame[post=s]{formula}, the target of the \sn{first-order semantics} of \Naproche,
      are recursively defined as follows:
      \begin{enumerate}
        \item $\defnotation\foltrue$ is a \sn{formula};
        \item $\defnotation\folfalse$ is a \sn{formula};
        \item $\defnotation\folthis$ is a \sn{formula};
        \item if $\phivar$ is a \sn{formula},
          then $\defnotation\folneg\phivar$ is a \sn{formula};
        \item if $\phivar$ and $\psivar$ are \sns{formula},
          then $\defnotation\folconj{\phivar,\psivar}$ is a \sn{formula};
        \item if $\phivar$ and $\psivar$ are \sns{formula},
          then $\defnotation\foldisj{\phivar,\psivar}$ is a \sn{formula};
        \item if $\phivar$ and $\psivar$ are \sns{formula},
          then $\defnotation\folimpl{\phivar}{\psivar}$ is a \sn{formula};
        \item if $\phivar$ and $\psivar$ are \sns{formula},
          then $\defnotation\folequiv{\phivar,\psivar}$ is a \sn{formula};
        \item if $\xvar$ is a \sn{variable declaration} and $\phivar$ a \sn{formula},
          then $\defnotation\folforall\xvar\phivar$ is a \sn{formula};
        \item if $\xvar$ is a \sn{variable declaration} and $\phivar$ a \sn{formula},
          then $\defnotation\folexists\xvar\phivar$ is a \sn{formula};
        \item if $\Tvar$ is a \sn{tag} and $\phivar$ a \sn{formula},
          then $\defnotation\foltag\Tvar\phivar$ is a \sn{formula};
        \item if $\ivar\in\mathbb N$,
          then $\defnotation\foldbindex\ivar$ is a \sn{formula},
          commonly called a \definame{De Brujin index};
        \item if $\vvar$ is a \sn{variable name},
          then $\vvar$ is a \sn{formula};
        \item if $\tvar!$ is a \sn{term name} and $\phiseq!$ are \sns{formula},
          then $\tvar\phiseq$ is a \sn{formula};
          in this case $\tvar\phiseq$ is called a \definame{term}.
      \end{enumerate}
    \end{sparagraph}
  \end{sfragment}
\end{sfragment}

\ifinputref\else\printlicense[CcByNcSa]{Marcel Schütz (2026)}\fi

\end{smodule}
\end{document}
